\section{Towards an NPE-based Shear Estimator}
\label{sec:shear}
\XW{Same model and test set as Sec.~\ref{sec:results} (100k training set, Xu3 only).}

In Sec.~\ref{sec:results}, we showed that our NPE produces well-calibrated posteriors. However, what we ultimately need is an estimate of the shear that is unbiased over an ensemble of galaxies, i.e. one with small $m$ and $\mathbf{c}$ in Eq.~(\ref{eqn:linear-bias}). In this section, we explore how the NPE posteriors can be turned into such an estimate. We start from the naive choice, the posterior mean of each galaxy, and show that it is strongly biased (Sec.~\ref{ssec:raw-bias}). We explain where this bias comes from in Sec.~\ref{ssec:shrinkage}, and then explore several ways to reduce it: weighting galaxies and cutting on their posterior widths (Sec.~\ref{ssec:weights-cuts}), calibrating the shear response (Sec.~\ref{ssec:response}), and, most importantly, combining the full likelihoods of galaxies that share the same shear instead of using point estimates (Sec.~\ref{ssec:stacking}). We present these as explorations rather than as a final, calibrated shear estimator framework, which we leave to future work.

\subsection{Raw Posterior Means}
\label{ssec:raw-bias}

For each test galaxy $i$, we take the mean of its TF-weighted posterior, $\hat{g}_{k,i}$, as a point estimate of its shear. We then measure $m$ and $c$ separately for the two shear components with a linear least-squares fit of $\hat{g}_{k,i} - g_{k,i} = m_k\, g_{k,i} + c_k$, which is Eq.~(\ref{eqn:linear-bias}) with $\hat{g}$ in place of $\hat{\epsilon}$. We use all Xu3 galaxies with equal weights and no cuts, and quote the standard errors of the fit as our uncertainties. We restrict the fit to galaxies with $|g_k| < 0.02$, since this range covers the shear of the vast majority of galaxies in cosmic shear measurements. At larger shears, the response of our estimator becomes slightly nonlinear, with a term cubic in the shear. Such third-order shear biases have also been found for other shear estimators, and become important at the large shears near galaxy clusters \citep{Anacal, AnaCalCluster}.

\begin{figure}
    \includegraphics[width=\linewidth]{figs/fig_raw_bias.pdf}
    \caption{Bias of the raw posterior means on the Xu3 test set ($10^5$ galaxies, equal weights, no cuts). Points show the mean of $\hat{g}_k - g_k$ in bins of the true shear for $g_1$ (blue circles) and $g_2$ (red squares); error bars are smaller than the markers. Solid lines are the linear fits over $|g_k| < 0.02$ (shaded), and dashed lines extend them to the full prior range.}
    \label{fig:raw-bias}
\end{figure}

Fig.~\ref{fig:raw-bias} shows the mean of $\hat{g}_k - g_k$ in bins of the true shear. The additive bias is consistent with zero for both components, but the multiplicative bias is large, with $m_1 = -0.616 \pm 0.014$ and $m_2 = -0.592 \pm 0.014$. In other words, the posterior means respond to only about 40\% of the true shear. Multiplicative biases of tens of percent are not unusual for uncalibrated shear estimators. In DES Y1, for example, the uncalibrated bias of \textsc{im3shape} ranges from $m\approx-0.1$ for high-SNR galaxies to $m\approx-0.3$ for the lowest-SNR galaxies \citep{DESY1Shape}, with similar levels found for HSC \citep{HSCY1Sims}. These biases are then calibrated with image simulations, which motivates the calibration we explore in Sec.~\ref{ssec:response}. Finding a significant multiplicative bias in the point estimate also does not contradict the calibration results of Sec.~\ref{ssec:calibration}. As we show in Sec.~\ref{ssec:shrinkage}, a bias of this kind is expected for posterior means when the data constrain the shear only weakly compared to the prior, even if the posteriors are correct.

We also measure the shape noise of each galaxy following \citep{KLDESI}, as $\sigma_{\epsilon,i} = \sqrt{\sigma_{g_1,i}^2 + \sigma_{g_2,i}^2}$, where $\sigma_{g_k,i}$ is the standard deviation of the TF-weighted posterior of $g_k$. Across the Xu3 test set, the median $\sigma_{\epsilon,i}$ is 0.066, with 68\% of galaxies between 0.040 and 0.079. About a third of the galaxies have $\sigma_{\epsilon,i} > 0.075$, close to the value of $\sqrt{2}\sigma_\pi = 0.082$ that we would get if the posteriors of both components were equal to the prior. Here $\sigma_\pi = 0.2/\sqrt{12} = 0.058$ is the standard deviation of the uniform prior $\mathcal{U}(-0.1, 0.1)$ on each shear component, and the factor of $\sqrt{2}$ comes from adding the two components in quadrature. The data of these galaxies carry little information about the shear, and as we will see in the next section, they are the main source of the bias in the posterior means. We compare our shape noise with that of \citep{KLDESI} in Sec.~\ref{sec:discussion}.

\subsection{Shrinkage of Posterior Means}
\label{ssec:shrinkage}

The bias of the posterior means can be understood as follows. Consider a single shear component $g$ of one galaxy, and suppose that its data give a Gaussian likelihood of width $\sigma_L$, centered on a noisy measurement $g_\mathrm{obs}$ whose average over noise realizations is the true shear. If the prior on $g$ is also a Gaussian centered on zero with width $\sigma_\pi$, the posterior is a Gaussian with mean and variance
\begin{equation}
\label{eqn:gauss-posterior}
    \hat{g} = \frac{\sigma_\pi^2}{\sigma_\pi^2 + \sigma_L^2}\, g_\mathrm{obs}, \qquad
    \sigma_\mathrm{post}^2 = \frac{\sigma_L^2\, \sigma_\pi^2}{\sigma_L^2 + \sigma_\pi^2}.
\end{equation}
The posterior mean is pulled toward the center of the prior by an amount that depends on how informative the data is compared to the prior. Averaging over noise realizations at a fixed true shear $g$ gives $\langle \hat{g} \rangle = g\,\sigma_\pi^2/(\sigma_\pi^2 + \sigma_L^2)$, so using Eq.~(\ref{eqn:gauss-posterior}) to write the result in terms of the posterior width, the multiplicative bias is
\begin{equation}
\label{eqn:shrinkage}
    m = -\frac{\sigma_L^2}{\sigma_L^2 + \sigma_\pi^2} = -\frac{\sigma_\mathrm{post}^2}{\sigma_\pi^2}.
\end{equation}
A galaxy with $\sigma_L \ll \sigma_\pi$ has $m \approx 0$, while a galaxy whose data carry no shear information has $\sigma_\mathrm{post} \to \sigma_\pi$ and $m \to -1$. For a sample of galaxies, the fitted $m$ is approximately the average of the values of the individual galaxies, $m \approx -\langle \sigma_\mathrm{post}^2 \rangle / \sigma_\pi^2$. Our shear prior is uniform rather than Gaussian, so we approximate it by a Gaussian with the same variance, $\sigma_\pi^2 = 0.2^2/12$.

This bias is a natural result from the posterior mean, and would be present even for exact posteriors. For correct posteriors, the average true shear of galaxies with a given $\hat{g}$ is equal to $\hat{g}$, but the average $\hat{g}$ of galaxies with a given true shear, which is what $m$ measures, is shrunk toward zero.

To test whether shrinkage explains the bias we measured in Sec.~\ref{ssec:raw-bias}, we compare Eq.~(\ref{eqn:shrinkage}) with the $m$ measured in groups of galaxies with different posterior widths. We split the test set into 15 equal-size bins of $\sigma_{\epsilon,i}$, and in each bin fit $m_k$ as in Sec.~\ref{ssec:raw-bias} and compute the average posterior variance $\langle \sigma_{\mathrm{post},k}^2 \rangle$. Note that grouping galaxies by their measured posterior widths is not the most stringent version of this test. The posterior width of a galaxy depends on its measured shear, with galaxies that appear more strongly sheared tending to have narrower posteriors. Selecting galaxies with narrow posteriors therefore preferentially selects galaxies whose measured shear $|\hat{g}|$ is larger than their true shear. This selection effect is strongest in the bins with the narrowest posteriors.

\begin{figure}
    \includegraphics[width=\linewidth]{figs/fig_shrinkage.pdf}
    \caption{Multiplicative bias of the raw posterior means against the average posterior variance, in 15 equal-size bins of $\sigma_{\epsilon,i}$ on the Xu3 test set (6,667 galaxies per bin), for $g_1$ (blue circles) and $g_2$ (red squares). The dashed line is the shrinkage prediction of Eq.~(\ref{eqn:shrinkage}). The dotted vertical line marks the prior variance $\sigma_\pi^2$.}
    \label{fig:shrinkage}
\end{figure}

Fig.~\ref{fig:shrinkage} shows the result. The measured $m$ follows the prediction from the bins with intermediate posterior widths to $m \approx -1$ in the bins with the widest posteriors. Over the whole test set, the prediction gives $m_1 = -0.607$ and $m_2 = -0.617$, close to the measured $m_1 = -0.616 \pm 0.014$ and $m_2 = -0.592 \pm 0.014$. Shrinkage therefore explains almost all of the bias of the raw posterior means, which is dominated by galaxies with wide posteriors. This bias can be reduced by giving less weight to these galaxies (Sec.~\ref{ssec:weights-cuts}) or by calibrating the response of the estimator (Sec.~\ref{ssec:response}). A more natural approach, however, is to not reduce each posterior to a point estimate at all, and instead to combine the likelihoods of all galaxies that share the same shear, utilizing the full information in the posteriors (Sec.~\ref{ssec:stacking}).
